
\subsubsection{2.1 Robustification Methods Requiring Group Annotations}

The standard approach to improving worst-group accuracy employs explicit group labels during training or fine-tuning. Distributionally Robust Optimization (DRO) [Ben-Tal et al., 2013] minimizes worst-case expected loss over a family of distributions. GroupDRO \citep{Sagawaetal2020} instantiates this objective by weighting training samples according to group membership, achieving approximately 90\% worst-group accuracy on Waterbirds when group labels are available at training time. Deep Feature Reweighting (DFR) \citep{Izmailovetal2022} takes a two-stage approach: an ERM model is trained to convergence, then only the final linear layer is retrained on a small group-balanced labeled validation set. DFR achieves approximately 97\% worst-group accuracy on Waterbirds and approximately 92\% on CelebA. One implication of DFR's success is that invariant features are encoded in pretrained ERM representations even when the linear classification head has learned to rely on spurious attributes. These methods require group annotations, which may be costly or unavailable in practice.

\subsubsection{2.2 Annotation-Free Robustification Methods}

Several methods approximate group labels without explicit annotation. Just Train Twice (JTT) \citep{Liuetal2021} trains an initial ERM model, identifies its misclassified training samples, and upweights those examples in a second training pass. The assumption is that minority-group samples are disproportionately misclassified. JTT achieves approximately 82–86\% worst-group accuracy on Waterbirds without group labels. Last-Layer Feature Reweighting (LFR) \citep{Ghaznavietal2023} uses loss-based resampling to approximate group-balanced training. These methods use the training loss as a proxy for group membership.

GEORGE \citep{Sohonietal2020} takes a representation-based approach: ERM penultimate-layer embeddings at training convergence are clustered to discover spurious groups, and the resulting cluster assignments are used as pseudo-labels for DRO. PruSC \citep{Kimetal2024} applies clustering to identify and prune spurious shortcut neurons. AGRO \citep{Paranjapeetal2022} uses adversarial discovery of error-prone groups followed by DRO. These methods cluster learned representations to derive group structure without annotation.

None of these methods explicitly characterize the conditions under which the clustering step succeeds. The configuration details under which published clustering results were obtained — in particular, the training epoch at which embeddings are extracted and the number of clusters used — are not fully specified in the cited works, making it unclear whether reported detection results transfer to the epoch-5, k=2 setting used by intervention methods.

\subsubsection{2.3 Shortcut Learning Mechanisms and Feature Representations}

Shah et al. [2020] establish that SGD exhibits a simplicity bias: it preferentially acquires high signal-to-noise ratio features early in training. In spurious correlation settings, this corresponds to early acquisition of spurious attributes. Geirhos et al. [2020] document shortcut learning across multiple modalities. Izmailov et al. [2022] demonstrate that ERM models richly encode invariant features in their representations even after training, which motivates representation-based analysis.

\subsubsection{2.4 Gap: Conditionality of Annotation-Free Detection}

The literature does not address the question of when clustering-based annotation-free detection produces reliable cluster-group alignment. The implicit assumption — that spurious features sufficiently strong to degrade worst-group accuracy are also strong enough to produce cluster-separable representations at epoch 5 — is not examined. The present work shows this assumption does not hold uniformly: it holds for scene-level spurious attributes (Waterbirds background habitat) but not for fine-grained texture-based spurious attributes (CelebA hair color) at the epoch-5, k=2 setting. The differentiating factor is the alignment between the spurious attribute and the ImageNet pretraining prior.

GEORGE \citep{Sohonietal2020} notes that clustering quality depends on representation quality but does not study the pretraining alignment factor. DFR \citep{Izmailovetal2022} demonstrates that ERM models encode invariant features in representations but does not study when clustering fails to recover spurious group structure. To the authors' knowledge, the feature-strength conditionality characterized in this work — that early-epoch detection reliability is determined by pretrained initialization alignment rather than by spurious correlation strength in the downstream task — has not been studied in prior work.

